% ============================================================
% Experiment 1 — REVISED (PURS distance, 5 000 users)
% Compare against: experiment1_original.tex (Jaccard, 100 users)
% ============================================================

\section{Experiment 1: Validating the Unimodal Serendipity Hypothesis}
\label{sec:experiment1}

\subsection{Motivation: The ``Goldilocks'' Question}

A fundamental question in serendipity research is whether there exists an
optimal level of unexpectedness --- not too similar, not too different, but
``just right'' --- that maximises the chance of a serendipitous discovery.
This intuition underlies what we call the \textbf{unimodal hypothesis}:

\begin{quote}
\textit{User satisfaction with unexpected items follows an inverted-U
relationship with unexpectedness: satisfaction is low when items are too
familiar, peaks at intermediate unexpectedness, and declines again when items
become too unfamiliar.}
\end{quote}

\textbf{If this hypothesis were false} --- if no such peak exists --- then
single-objective optimisation of unexpectedness would be provably insufficient,
and serendipity would require explicitly balancing unexpectedness against
relevance through multi-objective methods.  Experiment 1 tests whether the
unimodal relationship exists, and if so, where the peak lies relative to the
observable distance range.

% ============================================================
\subsection{Experimental Setup}
% ============================================================

This experiment follows the default configuration (Table~\ref{tab:default_setup}).
Unlike later experiments, no recommendations are generated.  Instead, PURS
unexpectedness scores (Eq.~\ref{eq:unexpectedness}) are computed for each
user's \emph{test-set} tracks and paired with the user's actual rating,
yielding 114\,690 (user, track) pairs across 5\,000 users.  We then ask
whether average rating varies systematically with unexpectedness, and whether
the relationship forms an inverted-U shape.

% ============================================================
\subsection{Statistical Analysis}
% ============================================================

\subsubsection{Binning and Summary Statistics}

Test pairs were grouped into ten equal-width distance bins $[0.0, 0.1),
[0.1, 0.2), \ldots, [0.9, 1.0)$.  For each bin we computed the mean rating,
standard deviation, number of pairs, and number of unique users.

\subsubsection{Quadratic Regression}

To test the inverted-U shape formally, we fit a quadratic regression of mean
bin rating on bin midpoint:

\begin{equation}
    \bar{r}(d) = \beta_0 + \beta_1 d + \beta_2 d^2
    \label{eq:quadratic}
\end{equation}

An inverted-U requires $\beta_2 < 0$ (downward curvature) and a peak
$\delta^* = -\beta_1 / (2\beta_2)$ within $[0, 1]$.  Model fit is assessed
via $R^2$; we consider $R^2 > 0.40$ an adequate fit.

\subsubsection{ANOVA}

A one-way ANOVA across the ten distance bins tests whether mean ratings differ
significantly between unexpectedness levels ($\alpha = 0.05$).

% ============================================================
\subsection{Results}
% ============================================================

\subsubsection{Bin Statistics}

Table~\ref{tab:exp1_bins} shows the mean rating and sample size for each
distance bin.  A clear pattern emerges: ratings are lowest when tracks are
very close to a user's interest clusters ($d < 0.1$, mean = 1.31), rise
steadily through intermediate distances, peak at $d \in [0.8, 0.9)$ (mean =
1.72), and then decline at maximum distance $d \in [0.9, 1.0)$ (mean = 1.57).

\begin{table}[h]
\centering
\caption{Mean rating and sample size per PURS distance bin (5\,000 users,
         114\,690 pairs). Ratings are binarised: 1 = disliked, 5 = liked.
         The percentage liked column converts the mean rating to an
         interpretable satisfaction rate.}
\label{tab:exp1_bins}
\begin{tabular}{lccc}
\toprule
\textbf{Distance bin} & \textbf{Mean rating} & \textbf{\% liked} & \textbf{Pairs} \\
\midrule
{[0.0--0.1)} & 1.312 &  7.8\% & 33\,815 \\
{[0.1--0.2)} & 1.518 & 12.9\% & 16\,257 \\
{[0.2--0.3)} & 1.481 & 12.0\% & 11\,441 \\
{[0.3--0.4)} & 1.549 & 13.7\% &  8\,472 \\
{[0.4--0.5)} & 1.531 & 13.3\% &  7\,250 \\
{[0.5--0.6)} & 1.531 & 13.3\% &  6\,217 \\
{[0.6--0.7)} & 1.589 & 14.7\% &  5\,719 \\
{[0.7--0.8)} & 1.663 & 16.6\% &  5\,269 \\
{[0.8--0.9)} & \textbf{1.719} & \textbf{18.0\%} &  4\,321 \\
{[0.9--1.0)} & 1.569 & 14.2\% & 15\,929 \\
\bottomrule
\end{tabular}
\end{table}

\noindent
The percentage-liked column makes the magnitude concrete: satisfaction at the
peak ($d \approx 0.85$, 18.0\%) is more than double that for near-expected
tracks ($d < 0.1$, 7.8\%), and is notably higher than at the maximum-distance
bin (14.2\%).

\subsubsection{Quadratic Regression}

The fitted model is:

\begin{equation}
    \bar{r}(d) = 1.336 + 0.684\,d - 0.396\,d^2
    \qquad R^2 = 0.721
    \label{eq:exp1_fit}
\end{equation}

All three conditions for an inverted-U are satisfied:
\begin{itemize}
    \item $\beta_2 = -0.396 < 0$: the relationship curves downward (\checkmark)
    \item $\delta^* = -0.684 / (2 \times -0.396) = 0.863$: peak lies within $[0, 1]$ (\checkmark)
    \item $R^2 = 0.721$: the quadratic model explains 72\% of the variance in
          bin-level ratings, indicating a strong fit (\checkmark)
\end{itemize}

\subsubsection{ANOVA}

The one-way ANOVA yields $F(9, 114{,}680) = 110.16$, $p \approx 0$, confirming
that mean ratings differ significantly across distance bins.  The null
hypothesis of equal means is decisively rejected.

\subsubsection{Visual Summary}

\textbf{[ADD FIGURE --- exp1\_purs\_unimodal\_relationship.png]}

\begin{figure}[htbp]
\centering
\includegraphics[width=0.95\textwidth]{MScTemplate/figs/exp1_purs_unimodal_relationship.png}
\caption{PURS Experiment 1 results (5\,000 users).
         \emph{(a)} Mean rating per distance bin with standard errors and
         quadratic fit ($R^2 = 0.721$); the dashed vertical line marks
         $\delta^* = 0.863$.
         \emph{(b)} Mean rating per bin (bar chart).
         \emph{(c)} Sample size distribution showing concentration near
         $d = 0$ (very familiar tracks) and $d = 1$ (cold/unknown items).
         \emph{(d)} Percentage of liked tracks by bin, confirming the
         inverted-U in satisfaction terms.}
\label{fig:exp1_relationship}
\end{figure}

% ============================================================
\subsection{Interpretation}
% ============================================================

\subsubsection{The Unimodal Hypothesis is Supported}

The inverted-U relationship is clear and statistically robust.  Users rate
tracks that are moderately-to-highly unexpected ($d \approx 0.86$) more
favourably than tracks that are very familiar ($d < 0.1$) or maximally
unfamiliar ($d \approx 1.0$).  This validates the Goldilocks intuition for
music: some degree of unexpectedness is beneficial, but the most extreme
surprises are also the least satisfying.

\subsubsection{Why the Peak is Near --- But Not At --- Maximum Distance}

The optimal distance $\delta^* = 0.863$ is high, but the decline beyond it
is real and interpretable.  The $[0.9, 1.0)$ bin contains the largest share
of \emph{cold items} --- tracks with no SVD embedding --- which receive a
default distance of 1.0.  These items lack any learned preference signal:
the model has no basis for predicting whether the user would like them.  The
drop in ratings at $d \to 1$ therefore reflects the failure of \emph{relevance}
at the extreme, not a failure of unexpectedness.

\subsubsection{Why This Still Motivates Multi-Objective Optimisation}

One might ask: if higher unexpectedness generally correlates with higher
satisfaction, why not simply maximise unexpectedness?  Two reasons argue
against this:

\begin{enumerate}
    \item \textbf{The peak is not at 1.0.}  Satisfaction declines beyond
          $\delta^* = 0.863$.  Maximising distance overshoots the optimum and
          returns worse recommendations than targeting the neighbourhood of
          $\delta^*$.

    \item \textbf{Unexpectedness alone cannot identify which distant items are
          worth recommending.}  Within the high-distance region ($d > 0.7$),
          satisfaction varies considerably across items.  Some highly unexpected
          tracks are genuinely delightful; others are simply irrelevant.
          Distinguishing between them requires a relevance signal --- precisely
          what the collaborative filtering score provides.  A system that
          maximises distance without a relevance filter will return a random
          sample of unexpected items, most of which the user will not enjoy.
\end{enumerate}

\noindent
The unimodal relationship therefore does not remove the need for
multi-objective optimisation; it \emph{characterises} the space in which that
optimisation must operate.  The ideal recommendation sits in the high-distance
region near $\delta^*$ \emph{and} has a high predicted rating --- a conjunction
of two objectives that neither can guarantee alone.

% ============================================================
\subsection{Limitations}
% ============================================================

\paragraph{Single distance metric.}
Results are based on PURS latent-space distance (Euclidean in SVD embedding
space, $k = 3$ clusters).  Jaccard distance on style tags produced a weaker
signal ($R^2 = 0.37$, hypothesis not supported at the $R^2 > 0.40$ threshold),
suggesting the finding is sensitive to how unexpectedness is operationalised.
The methodology section discusses this choice in detail.

\paragraph{Binary ratings.}
AMBAR ratings were binarised (liked / disliked) to address the skewed
distribution of the original 1--5 scale.  This reduces rating granularity and
may compress differences between bins.

\paragraph{Threshold dependency.}
The $R^2 > 0.40$ acceptance criterion for the quadratic fit is a
methodological choice.  The underlying pattern (inverted-U with $\delta^* =
0.863$, $F = 110.16$, $p \approx 0$) is robust regardless of this threshold.

% ============================================================
\subsection{Conclusion}
% ============================================================

Experiment 1 confirms the unimodal serendipity hypothesis in the AMBAR music
domain.  Across 5\,000 users and 114\,690 test-set ratings, user satisfaction
with unexpected tracks follows a clear inverted-U relationship with
PURS latent-space distance ($R^2 = 0.721$, $\delta^* = 0.863$, $F = 110.16$,
$p \approx 0$).  Satisfaction is lowest for very familiar tracks, rises through
intermediate unexpectedness levels, peaks near the high end of the distance
range ($d \approx 0.86$), and declines for maximally unfamiliar items.

Critically, the peak at $\delta^* = 0.863$ does not occur automatically: it
emerges only when recommendations combine high unexpectedness with sufficient
relevance.  A system that blindly maximises unexpectedness will overshoot the
peak and return items that are surprising but unlikeable.  This establishes
the core motivation for Experiment 2: rather than optimising a single
objective, we explore weighted combinations of unexpectedness and relevance,
systematically searching for the parameter $\alpha \in [0, 1]$ --- and
ultimately a per-user vector $\boldsymbol{\alpha}$ --- that produces the
highest aggregate serendipity.
